% 46
\begin{frame}{Плохо обусловленная задача}
\small
Пусть
\[
\kappa=100.
\]
Тогда
\[
q=1-\frac1{100}=0.99.
\]
За одну итерацию гарантированно исчезает лишь около $1\%$ текущей ошибки.

\centering
\includegraphics[width=0.62\linewidth]{kappa_rates.pdf}
\end{frame}

% 47
\begin{frame}{Сколько итераций нужно в сильно выпуклом случае?}
\small
Хотим
\[
\left(1-\frac1\kappa\right)^k
\le
\varepsilon.
\]
Используем
\[
1-u\le e^{-u},
\qquad u\ge0.
\]
Тогда
\[
\left(1-\frac1\kappa\right)^k
\le
e^{-k/\kappa}.
\]
\end{frame}

% 48
\begin{frame}{Откуда возникает логарифм}
\small
Чтобы
\[
e^{-k/\kappa}\le\varepsilon,
\]
достаточно
\[
-\frac{k}{\kappa}\le\log\varepsilon.
\]
То есть
\[
\boxed{
k\ge\kappa\log\frac1\varepsilon.}
\]
В обозначениях порядка:
\[
\boxed{
k=O\!\left(\kappa\log\frac1\varepsilon\right).}
\]
\end{frame}

% 49
\begin{frame}{Сравнение двух классов функций}
\small
\centering
\begin{tabular}{lcc}
\toprule
& Выпуклая + гладкая & Сильно выпуклая + гладкая\\
\midrule
Ошибка & $O(1/k)$ & $O(q^k)$\\
До точности $\varepsilon$ & $O(1/\varepsilon)$ & $O(\log(1/\varepsilon))$\\
Минимум единственен & не обязательно & да\\
Роль $\kappa$ & не обязательна & определяет скорость\\
\bottomrule
\end{tabular}

\vspace{0.7em}
\includegraphics[width=0.53\linewidth]{complexity_epsilon.pdf}
\end{frame}

% 50
\begin{frame}{Связь с квадратичными задачами}
\small
Для
\[
f(x)=\frac12x^\top Ax-b^\top x,
\qquad A\succ0,
\]
имеем
\[
L=\lambda_{\max}(A),
\qquad
\mu=\lambda_{\min}(A).
\]
Поэтому
\[
\boxed{
\kappa=\frac{L}{\mu}
=\frac{\lambda_{\max}(A)}{\lambda_{\min}(A)}.
}
\]
Это ровно то же число обусловленности, которое появилось на лекции 3.
\end{frame}

% 51
\begin{frame}{Один минимум, разные числа обусловленности}
\centering
\includegraphics[width=0.96\linewidth]{gd_trajectories_kappa.pdf}

\vspace{0.1em}
\small
Чем более вытянуты линии уровня, тем дольше градиентный спуск избавляется от медленной компоненты ошибки.
\end{frame}

% 52
\begin{frame}{Фактическая ошибка и теоретическая граница}
\centering
\includegraphics[width=0.76\linewidth]{objective_gaps.pdf}

\vspace{0.15em}
\small
Пунктир — гарантированная скорость $\left(1-1/\kappa\right)^k$; сплошные линии — реальный эксперимент.
\end{frame}

% 53
\begin{frame}{Зачем используют логарифмический масштаб}
\small
Если
\[
e_k=Cq^k,
\qquad 0<q<1,
\]
то после логарифмирования
\[
\log e_k=\log C+k\log q.
\]
Это линейная функция от $k$.

\begin{block}{Практический признак}
Почти прямая линия ошибки на полулогарифмическом графике — характерный признак геометрической сходимости.
\end{block}
\end{frame}

% 54
\begin{frame}{Теория против эксперимента}
\centering
\includegraphics[width=0.73\linewidth]{theory_vs_experiment.pdf}

\vspace{0.15em}
\small
Теоретическая оценка строится для всего класса задач и потому часто консервативна. Она должна быть над фактической ошибкой, а не совпадать с ней.
\end{frame}

% 55
\begin{frame}{Самопроверка}
\small
Пусть
\[
L=20,
\qquad
\mu=2.
\]
Ответьте:
\begin{enumerate}
\item Чему равно $\kappa$?
\item Чему равен множитель $q=1-1/\kappa$?
\item Во сколько раз примерно уменьшится верхняя оценка ошибки после 20 шагов?
\end{enumerate}

\[
\boxed{\kappa=10,\qquad q=0.9,\qquad q^{20}\approx0.12.}
\]
\end{frame}

% 56
\begin{frame}{Карта лекции}
\small
Для выпуклой $L$-гладкой функции:
\[
\boxed{
\gap{x_k}
\le
\frac{L\norm{x_0-x^\star}^2}{2k}.
}
\]
Добавляем $\mu$-сильную выпуклость:
\[
\boxed{
\gap{x_k}
\le
\left(1-\frac\mu L\right)^k\bigl(f(x_0)-f^\star\bigr).
}
\]
Именно отношение
\[
\boxed{\kappa=L/\mu}
\]
управляет худшей гарантированной геометрической скоростью.
\end{frame}

% 57
\begin{frame}{Следующий естественный вопрос}
\small
При плохой обусловленности
\[
\kappa\gg1
\]
обычный градиентный спуск может быть медленным, даже если задача выпуклая и очень хорошо изучена.

\vspace{0.6em}
\begin{alertblock}{Что можно сделать?}
Можно ли изменить алгоритм так, чтобы быстрее проходить длинные узкие долины и меньше зависеть от $\kappa$?
\end{alertblock}

Это приведёт нас к более продвинутым методам оптимизации.
\end{frame}
